\section{دوبامين بومضة ضوء}
\label{sec:dishdopamine}

التجربة المباشرة الوحيدة التي نعرفها، حيث أُعطيت المزرعة الدوبامين معلّمًا، أُجريت على منصة عضيّات يتصل بها الباحثون عن بُعد~\cite{Jordan2024}. أُضيف إلى السائل الذي يغمر العضيّات دوبامين داخل ”قفص“ حساس للضوء: فالجزيء المقيَّد لا يؤثر في المستقبلات. وتركيز الدوبامين المقيَّد $1\mM$. وحين يلزم أن تُكافأ الشبكة يُسلَّط عليها ضوء فوق بنفسجي: فينفتح القفص ويخرج الدوبامين إلى الحجم.

والحلقة مبنية هكذا. يُقدَّم المنبه نفسه مرة بعد مرة؛ فإذا أثار نبضات أكثر من ذي قبل أُطلقت الومضة وتحرر الدوبامين. وعلى مدى $240$ تكرارًا ازداد عدد النبضات المستثارة إجمالًا. ويكتب الباحثون أنفسهم أنه لا يمكن من تجربة واحدة كهذه أن يُستنتج أن الزيادة سببها الدوبامين~\cite{Jordan2024}.

هذه، في عين المهندس، أول محاولة لتجميع قاعدة العوامل الثلاثة~\eqref{eq:threefactor} في طبق: فعالية المرسل والمستقبِل تترك أثرًا، والإشارة العامة تقرر هل يُثبَّت التغيير. لكن الإشارة هنا موجبة فقط: المكافأة إما موجودة وإما غائبة، والجهاز لا يُصدر خطأً سالبًا $\delta<0$. ثم إن الدوبامين ينتشر ويُزال ببطء، كالتركيز في~\eqref{eq:lowpass}، فقد ترفع الومضات المتكررة مستواه العام لا غير. ولتمييز التعلم من ذلك تلزم تجربتان ضابطتان: العدد نفسه من الومضات في لحظات عشوائية لا صلة لها باستجابة الشبكة، والومضات نفسها بلا دوبامين مقيَّد. ويلزم أيضًا اختبار بعد الراحة: هل بقي التغيير بعد أن زال الدوبامين؟

\section{الدوبامين يعلّم السيليكون}
\label{sec:biohybrid}

في التجربة المعاكسة تعلّم الخلايا الحية الإلكترونيات~\cite{Keene2020}. تُستنبت خلايا تفرز الدوبامين في قناة دقيقة فوق عنصر عصبي الشكل عضوي، أي ترانزستور تؤدي موصلية قناته دور وزن المشبك. والدوبامين الخارج من الخلايا يتأكسد عند قطب العنصر، فيغيّر ذلك الموصلية. ويحاكي الجهاز أيضًا آلية إزالة الناقل من الشق، فيمكن أن يُغيَّر الوزن تغييرًا دائمًا وأن يُعاد أيضًا إلى ما كان عليه.

من حيث الدارة هذا مكامل تسريبي بدخل كيميائي: الوزن يراكم تدفق الدوبامين، و”الإزالة“ تحدد سرعة نسيانه، وهي دارة \foreignlanguage{english}{RC} نفسها التي في~\eqref{eq:lowpass}. والمهم هنا اتجاه الاتصال. فمسبار الفصل~\ref{sec:debugging} لا يرى سجلات البثّ لأنها كيميائية. وهذا العنصر يقرأ سجلًا كيميائيًا بالذات ويحوّله إلى وزن كهربائي. وهو، بالنسبة إلى واجهات الدماغ والحاسوب، طريق إلى مستشعر يسمع لا ناقل البيانات وحده، بل سجلات التحكم أيضًا.

\section{عضيّات بعصبوناتها الدوبامينية}
\label{sec:midbrainorg}

صار في الإمكان أن تُستنبت من الخلايا الجذعية البشرية عضيّات تشبه منطقة الدماغ التي توجد فيها عصبونات \nt{DOPA}. وفيها عصبونات دوبامينية عاملة تفرز الدوبامين~\cite{Jo2016}. وتُستنبت هذه العضيّات في الغالب لدراسة الأمراض. ولم نجد تجارب استُعمل فيها دوبامينها معلّمًا لشبكة أخرى.

وهذه، بالنسبة إلى الآلة المصنوعة من عصبونات حية، هي القطعة الناقصة. فمنصة الفصل~\ref{sec:livingmachine} ناقل بيانات وتحكم في التدفق بلا سجلات تحكم؛ وهنا مصدر إشارة البثّ موجود. ما ينقص هو الناقد: شبكة تحسب خطأ التنبؤ~\eqref{eq:ctd} وتتحكم في هذه العصبونات. ووصل عضيّة قشرية بعضيّة كهذه بحيث يُفرز الدوبامين بحسب الخطأ مسألة مفتوحة، وهو حتى الآن فرضية لا تجربة.

\section{عضيّة توازن عمودًا بلا دوبامين}
\label{sec:cartpole}

أكمل التجارب الحديثة يستغني عن الدوبامين تمامًا~\cite{Robbins2026}. وُصلت عضيّات قشرة الفأر في حلقة مغلقة بمسألة ”العمود على العربة“: فعالية العضيّة تحرّك العربة، وموضع العمود يعود إليها بالتحفيز. وبين المحاولات تتلقى العضيّة منبهات تعليم قصيرة عالية التردد. أيّها بالضبط، يختاره برنامج تعلّم معزَّز.

النتائج مقيسة:
\begin{itemize}
\item في أغلب العضيّات أعطت إشارات التعليم التي اختارها البرنامج نتيجة أفضل من إشارات مختارة عشوائيًا أو من غياب الإشارة؛
\item بعد $45$ دقيقة من الراحة لم يبقَ التحسن؛
\item حجب مستقبلات الغلوتامات بنوعيها، \foreignlanguage{english}{AMPA} و\foreignlanguage{english}{NMDA}، ألغى التحسن.
\end{itemize}

البند الأخير يدل على مكان ما تعلّمته العضيّة: في مشابك \nt{GLUT}، وعبر المستقبِل الذي يعمل في القسم~\ref{sec:weightstore} كاشفًا لتزامن الدخل والخرج. والبند الثاني يبيّن ما ينقص: التغير السريع موجود، والتثبيت غائب، كالمتعلم السريع في الفصل~\ref{sec:motorlearning} بلا المتعلم البطيء. أما دور المعلّم فتغيّر: حلّ برنامج محل المهندس الذي يختار نمط التيار، لكن المعلّم لا يزال خارج الطبق.

ولم نجد كذلك بين التجارب الأقدم على تعليم المزارع فوق مصفوفات الأقطاب تجربة واحدة استُعمل فيها الدوبامين: إشارة التعليم في كل مكان كانت تحفيزًا كهربائيًا.

\section{من يعلّم المزرعة}

\begin{table}[t]
\centering
\footnotesize
\setlength{\tabcolsep}{4pt}
\renewcommand{\arraystretch}{1.15}
\begin{tabular}{>{\raggedright}p{2.7cm}>{\raggedright}p{2.5cm}>{\raggedright}p{3.2cm}>{\raggedright\arraybackslash}p{4.4cm}}
\toprule
التجربة & مَن يعلّم & إشارة التعليم & ما هو معروف \\
\midrule
إيقاف التحفيز عند الاستجابة المطلوبة & المهندس & انتهاء التحفيز & تثبت الاستجابة، وهذا مُقاس \\
\foreignlanguage{english}{DishBrain} (الفصل~\ref{sec:dishbrain}) & المهندس & تيار متوقَّع أو عشوائي & نمو ذو دلالة في الجولات داخل الجلسة مع التغذية الراجعة الكاملة فقط، وأثر أصغر مع الصمت، وهذا مُقاس؛ غير ثابت من جلسة إلى أخرى؛ والسبب موضع نقاش \\
العمود على العربة & برنامج تعلّم معزَّز & منبهات عالية التردد يختارها البرنامج & أفضل من العشوائية، لكنه لا يبقى بعد الراحة، وهذا مُقاس \\
دوبامين بالومضة & قاعدة ”نبضات أكثر = مكافأة“ & دوبامين يحرّره الضوء & زيادة النبضات ملاحظة؛ دور الدوبامين غير مثبت \\
مشبك هجين حيوي & خلايا حية & دوبامين من الخلايا & تغيّر طويل لوزن العنصر وإعادته، وهذا مُقاس \\
عضيّات بعصبونات \nt{DOPA} & --- & --- & مصدر الدوبامين موجود؛ ولم يُستعمل معلّمًا \\
\bottomrule
\end{tabular}
\caption{من يعلّم العصبونات الحية على الشريحة، وبماذا.}
\label{tab:dishteachers}
\end{table}

في كل التجارب التي تتعلم فيها المزرعة نفسها، المعلّم حتى الآن في الخارج (الجدول~\ref{tab:dishteachers}): مهندس، أو برنامج، أو قاعدة وضعها مهندس. ولم يدخل الدوبامين الطبق إلا مرة واحدة، ودوره لم يُثبت بعد. أما تجميع قناة بثّ حقيقية في الطبق، بناقد وخطأ وأثر كما في الفصل~\ref{sec:dofa}، فهو الخطوة التالية التي لم يخطُها أحد بعد.
